% ============================================================
% UGA-Themed Beamer Presentation
% Literature Review: Love of Variety
%
% Compile with:
%   pdflatex literature_pres.tex
% ============================================================

\documentclass[aspectratio=169,professionalfonts]{beamer}

% ---------- COLORS ----------
\usepackage{xcolor}
\definecolor{UGARed}{HTML}{BA0C2F}
\definecolor{UGABlack}{HTML}{000000}
\definecolor{UGAGray}{HTML}{8A8D8F}

% ---------- THEME ----------
\usetheme{default}
\useinnertheme{circles}
\useoutertheme[subsection=false]{miniframes}
\usefonttheme{professionalfonts}
\setbeamertemplate{navigation symbols}{}

% Head/foot colors
\setbeamercolor{subsection in head/foot}{fg=white, bg=UGARed}
\setbeamercolor{headline}{fg=white, bg=UGARed}
\setbeamercolor{section in head/foot}{fg=white, bg=UGARed}
\setbeamercolor{section in head/foot shaded}{fg=white, bg=white}
\setbeamercolor{footline}{bg=UGABlack, fg=white}
\setbeamercolor{frametitle}{fg=UGABlack}
\setbeamercolor{title}{fg=UGABlack}
\setbeamercolor{normal text}{fg=UGABlack}
\setbeamercolor{structure}{fg=UGARed}

% Headline
\setbeamertemplate{headline}{
  \begin{beamercolorbox}[wd=\paperwidth,ht=2.5ex,dp=1.125ex]{headline}
    \insertsectionnavigationhorizontal{\paperwidth}{}{}
  \end{beamercolorbox}
}

% Footline
\setbeamertemplate{footline}{%
  \leavevmode%
  \hbox{%
    \begin{beamercolorbox}[wd=\paperwidth,ht=2.5ex,dp=1ex,leftskip=1em,rightskip=1em]{footline}%
      \ifnum\theframenumber=0\relax\else%
        \insertshortauthor{} \hfill \insertframenumber/\inserttotalframenumber%
      \fi%
    \end{beamercolorbox}}%
}

% ---------- PACKAGES ----------
\usepackage{appendixnumberbeamer}
\usepackage{amsmath, amssymb, mathtools}
\usepackage{graphicx}
\usepackage{booktabs}
\usepackage{array}
\usepackage{hyperref}
\usepackage{tcolorbox}

% ---------- METADATA ----------
\title[Love of Variety: Lit.\ Review]{Literature Review: Love of Variety}
\author[Tate Mason]{Tate Mason}
\institute[UGA]{John Munro Godfrey Sr.\ Dept.\ of Economics \\ University of Georgia}
\date{April 2026}

% ---------- AUTO-AGENDA ----------
\AtBeginSection[]{
  \begin{frame}[noframenumbering]{Roadmap}
    \tableofcontents[currentsection]
  \end{frame}
}

% ---------- MACROS ----------
\newcommand{\hl}[1]{\textbf{\textcolor{UGARed}{#1}}}
\hypersetup{colorlinks=true, linkcolor=UGARed, urlcolor=UGARed}

% ============================================================
\begin{document}

% ------------------------------------------------------------
% Title
% ------------------------------------------------------------
\begin{frame}[noframenumbering]
  \titlepage
\end{frame}

% ============================================================
\section{Motivation}
% ============================================================

\begin{frame}{Motivation}
  \begin{columns}[T]
    \begin{column}{0.55\textwidth}
      \hl{The problem with standard logit:}
      \begin{itemize}
        \item A new product always raises utility via a fresh $\varepsilon$ draw
        \item No actual parameter for variety-seeking
        \item Consumers seem to prefer goods that are \emph{novel but familiar}
      \end{itemize}
      \vspace{6pt}
      \hl{This project models:}
      \begin{itemize}
        \item How deviation from a consumer's \emph{experienced mean} affects utility
        \item How advertising interacts with variety 
      \end{itemize}
    \end{column}
    \begin{column}{0.42\textwidth}
      \begin{tcolorbox}[colback=red!5!white, colframe=UGARed, title=\textbf{Three Literatures}]
        \begin{enumerate}
          \item Variety-seeking \& habit formation
          \item Dyanmic discrete-choice demand estimation
          \item Dynamic advertising
        \end{enumerate}
      \end{tcolorbox}
      \vspace{1pt}
      \begin{tcolorbox}[colback=red!5!white, colframe=UGARed, title=\textbf{Data}]
        Nielsen-Kilts Homescan and Marketscan scanner data 
      \end{tcolorbox}
    \end{column}
  \end{columns}
\end{frame}

% ============================================================
\section{Variety-Seeking \& Habit Formation}
% ============================================================

\begin{frame}{Dixit \& Stiglitz (1977): Love of Variety}
  \begin{tcolorbox}[colback=red!5!white, colframe=UGARed, title=\textbf{Citation}]
    Dixit, A. K. and Stiglitz, J. E. (1977). Monopolistic competition and optimum
    product diversity. \textit{American Economic Review}, 67(3), 297--308.
  \end{tcolorbox}
  \vspace{6pt}
  \begin{columns}[T]
    \begin{column}{0.55\textwidth}
      \hl{What they do:}
      \begin{itemize}
        \item Model preferences over a continuum of differentiated products
        \item Show consumers prefer greater variety under monopolistic competition
        \item Derive optimal product diversity conditions
      \end{itemize}
    \end{column}
    \begin{column}{0.42\textwidth}
      \hl{Relevance:} The canonical ``love of variety'' framework
    \end{column}
  \end{columns}
\end{frame}

\begin{frame}{McAllister \& Pessemier (1982): Variety-Seeking Review}
  \begin{tcolorbox}[colback=red!5!white, colframe=UGARed, title=\textbf{Citation}]
    McAllister, I. and Pessemier, E. A. (1982). Variety seeking behavior: An
    interdisciplinary review. \textit{Journal of Consumer Research}, 9(3), 311--322.
  \end{tcolorbox}
  \vspace{6pt}
  \begin{columns}[T]
    \begin{column}{0.55\textwidth}
      \hl{What they do:}
      \begin{itemize}
        \item Comprehensive review of variety-seeking literature through early 1980s
        \item Synthesize psychological, sociological, and economic perspectives
        \item Catalog empirical patterns in variety-seeking behavior
      \end{itemize}
    \end{column}
    \begin{column}{0.42\textwidth}
      \hl{Relevance:} Foundational reference for understanding motivations behind
      variety-seeking; useful for grounding the model's behavioral assumptions
    \end{column}
  \end{columns}
\end{frame}

\begin{frame}{Becker \& Murphy (1988): Rational Addiction}
  \begin{tcolorbox}[colback=red!5!white, colframe=UGARed, title=\textbf{Citation}]
    Becker, G. S. and Murphy, K. M. (1988). A theory of rational addiction.
    \textit{Journal of Political Economy}, 96(4), 675--700.
  \end{tcolorbox}
  \vspace{6pt}
  \begin{columns}[T]
    \begin{column}{0.55\textwidth}
      \hl{What they do:}
      \begin{itemize}
        \item Model addictive consumption as rational forward-looking behavior
        \item Past consumption raises the marginal utility of current consumption
        \item Embed habit stock dynamics into utility maximization
      \end{itemize}
    \end{column}
    \begin{column}{0.42\textwidth}
      \hl{Relevance:} Habit-stock approach maps onto the experienced-mean state variable
      $\bar{X}_{it}$; provides microfoundations for why $\bar{X}_{it}$ is welfare-relevant
    \end{column}
  \end{columns}
\end{frame}

\begin{frame}{Bronnenberg, Dub\'e \& Gentzkow (2021): Brand Preferences}
  \begin{tcolorbox}[colback=red!5!white, colframe=UGARed, title=\textbf{Citation}]
    Bronnenberg, B. J., Dub\'e, J.-P., and Gentzkow, M. (2021). Millenials and the Take-off of Craft Brands: Preference Formation in the U.S. Beer Industry. \\
    \textit{Working Paper}.
  \end{tcolorbox}
  \vspace{6pt}
  \begin{columns}[T]
    \begin{column}{0.55\textwidth}
      \hl{What they do:}
      \begin{itemize}
        \item Use consumer migration across markets as identification strategy
        \item Examine craft beer industry; find strong persistence in brand preferences
        \item Construct a ``capital stock'' measure to track consumption history
      \end{itemize}
    \end{column}
    \begin{column}{0.42\textwidth}
      \hl{Relevance:} Capital stock $\approx$ this project's $\bar{X}_{it}$; inspires
      the ``prior mean'' initialization in simulations. Key empirical motivation for persistence
    \end{column}
  \end{columns}
\end{frame}

% ============================================================
\section{Discrete-Choice Demand Estimation}
% ============================================================

\begin{frame}{McFadden (1974): Conditional Logit}
  \begin{tcolorbox}[colback=red!5!white, colframe=UGARed, title=\textbf{Citation}]
    McFadden, D. (1974). Conditional logit analysis of qualitative choice behavior.
    In Zarembka, P. (ed.), \textit{Frontiers in Econometrics}. Academic Press.
  \end{tcolorbox}
  \vspace{6pt}
  \begin{columns}[T]
    \begin{column}{0.55\textwidth}
      \hl{What they do:}
      \begin{itemize}
        \item Derive the multinomial logit from random utility maximization
        \item Show IIA property and its implications
        \item Establish the statistical framework for discrete choice
      \end{itemize}
    \end{column}
    \begin{column}{0.42\textwidth}
      \hl{Relevance:} The MNL is the baseline for this project's choice probabilities.
      The $\varepsilon_{ijt} \sim \text{Gumbel}$ assumption driving $P_{ijt}$ comes
      directly from this framework
    \end{column}
  \end{columns}
\end{frame}

\begin{frame}{Berry, Levinsohn \& Pakes (1995): BLP}
  \begin{tcolorbox}[colback=red!5!white, colframe=UGARed, title=\textbf{Citation}]
    Berry, S., Levinsohn, J., and Pakes, A. (1995). Automobile prices in market
    equilibrium. \textit{Econometrica}, 63(4), 841--890.
  \end{tcolorbox}
  \vspace{6pt}
  \begin{columns}[T]
    \begin{column}{0.55\textwidth}
      \hl{What they do:}
      \begin{itemize}
        \item Introduce random-coefficients logit for differentiated products
        \item Solve the market-share inversion problem via contraction mapping
        \item Provide IV strategy to handle price endogeneity
      \end{itemize}
    \end{column}
    \begin{column}{0.42\textwidth}
      \hl{Relevance:} The estimation framework for this project. BLP instruments will
      identify preferences for sameness and variety from observed market shares; random coefficients
      accommodate heterogeneity in variety preferences
    \end{column}
  \end{columns}
\end{frame}

\begin{frame}{Nevo (2001): BLP Applied to Cereal}
  \begin{tcolorbox}[colback=red!5!white, colframe=UGARed, title=\textbf{Citation}]
    Nevo, A. (2001). Measuring market power in the ready-to-eat cereal industry.
    \textit{Econometrica}, 69(2), 307--342.
  \end{tcolorbox}
  \vspace{6pt}
  \begin{columns}[T]
    \begin{column}{0.55\textwidth}
      \hl{What they do:}
      \begin{itemize}
        \item Apply BLP to the ready-to-eat cereal industry
        \item Carefully construct instruments; demonstrate identification
        \item Recover markups and market power estimates
      \end{itemize}
    \end{column}
    \begin{column}{0.42\textwidth}
      \hl{Relevance:} Practical implementation guide for BLP; cereal is a close cousin
      to the CPG scanner data setting. Blueprint for the estimation phase of this project
    \end{column}
  \end{columns}
\end{frame}

\begin{frame}{Conlon \& Gortmaker (2020): PyBLP Best Practices}
  \begin{tcolorbox}[colback=red!5!white, colframe=UGARed, title=\textbf{Citation}]
    Conlon, C. and Gortmaker, J. (2020). Best practices for differentiated products demand
    estimation with PyBLP. \textit{RAND Journal of Economics}, 51(4), 1108--1161.
  \end{tcolorbox}
  \vspace{6pt}
  \begin{columns}[T]
    \begin{column}{0.55\textwidth}
      \hl{What they do:}
      \begin{itemize}
        \item Provide numerical best practices for BLP-style estimators
        \item Document common pitfalls (integration rules, convergence tolerance)
        \item Release PyBLP, an open-source Python implementation
      \end{itemize}
    \end{column}
    \begin{column}{0.42\textwidth}
      \hl{Relevance:} Directly relevant for the structural estimation phase. PyBLP will
      be the estimation engine; this paper gives guidance on numerical implementation choices
    \end{column}
  \end{columns}
\end{frame}

\begin{frame}{Wei \& Jiang (2024): Neural Network Estimation}
  \begin{tcolorbox}[colback=red!5!white, colframe=UGARed, title=\textbf{Citation}]
    Wei, Y. and Jiang, Z. (2024). Estimating parameters of structural models using
    neural networks. \textit{Marketing Science}, 44(1), 102--128.
  \end{tcolorbox}
  \vspace{6pt}
  \begin{columns}[T]
    \begin{column}{0.55\textwidth}
      \hl{What they do:}
      \begin{itemize}
        \item Implement neural networks as a flexible moment-matching estimator
        \item Demonstrate efficiency gains over classical GMM in structural models
        \item Show improved accuracy given correct moment selection
      \end{itemize}
    \end{column}
    \begin{column}{0.42\textwidth}
      \hl{Relevance:} Potential computational shortcut for estimating the dynamic model
      on Nielsen-Kilts data; ties into Song (2025) as a practical implementation strategy
    \end{column}
  \end{columns}
\end{frame}

% ============================================================
\section{Dynamic Advertising}
% ============================================================

\begin{frame}{Nerlove \& Arrow (1962): Goodwill}
  \begin{tcolorbox}[colback=red!5!white, colframe=UGARed, title=\textbf{Citation}]
    Nerlove, M. and Arrow, K. J. (1962). Optimal advertising policy under dynamic
    conditions. \textit{Economica}, 29(114), 129--142.
  \end{tcolorbox}
  \vspace{6pt}
  \begin{columns}[T]
    \begin{column}{0.55\textwidth}
      \hl{What they do:}
      \begin{itemize}
        \item Model advertising as a stock of ``goodwill'' that depreciates over time
        \item Derive the firm's optimal advertising path as a dynamic control problem
        \item Show advertising investment trades off current cost vs.\ future demand
      \end{itemize}
    \end{column}
    \begin{column}{0.42\textwidth}
      \hl{Relevance:} Theoretical foundation for the firm's dynamic advertising problem.
      The goodwill stock maps onto consumer familiarity, linking naturally to the
      habit formation side
    \end{column}
  \end{columns}
\end{frame}

\begin{frame}{Erdem \& Keane (1996): Advertising as Learning}
  \begin{tcolorbox}[colback=red!5!white, colframe=UGARed, title=\textbf{Citation}]
    Erdem, T. and Keane, M. P. (1996). Decision-making under uncertainty: Capturing
    dynamic brand choice processes in turbulent consumer goods markets.
    \textit{Marketing Science}, 15(1), 1--20.
  \end{tcolorbox}
  \vspace{6pt}
  \begin{columns}[T]
    \begin{column}{0.55\textwidth}
      \hl{What they do:}
      \begin{itemize}
        \item Model advertising as a signal that updates consumer beliefs about quality
        \item Consumers Bayesian-update perceptions from ad exposure
        \item Estimate dynamic brand choice model on detergent data
      \end{itemize}
    \end{column}
    \begin{column}{0.42\textwidth}
      \hl{Relevance:} Estimation procedure and ``perceptions formed via advertising'' are
      directly relevant. This project \emph{flips} the direction: consumer choices inform
      firm advertising strategy
    \end{column}
  \end{columns}
\end{frame}

\begin{frame}{Dub\'e, Hitsch \& Rossi (2005): Adverising Timing}
  \begin{tcolorbox}[colback=red!5!white, colframe=UGARed, title=\textbf{Citation}]
    Dub\'e, J.-P., Hitsch, G. J., and Rossi, P. E. (2005).  An Empirical Model of Adverising Dynamics.
    \textit{Quantitative Marketing and Economics}, 3, 107-144.
  \end{tcolorbox}
  \vspace{6pt}
  \begin{columns}[T]
    \begin{column}{0.55\textwidth}
      \hl{What they do:}
      \begin{itemize}
        \item Utilize goodwill stock with distributed lag to model advertising as investment process
        \item Solve for Markov perfect equilibrium in a dynamic oligopoly game
        \item Using frozen entree data, find that their model overpredicts advertising persistence
      \end{itemize}
    \end{column}
    \begin{column}{0.42\textwidth}
      \hl{Relevance:} Update on the Nerlove-Arrow framework; the distributed lag structure informs the specification of the firm's advertising cost function. 
    \end{column}
  \end{columns}
\end{frame}

\begin{frame}{Bergemann \& Bonatti (2011): Optimal Targeting}
  \begin{tcolorbox}[colback=red!5!white, colframe=UGARed, title=\textbf{Citation}]
    Bergemann, D. and Bonatti, A. (2011). Targeting in advertising markets: Implications
    for offline vs.\ online media. \textit{RAND Journal of Economics}, 42(3), 417--443.
  \end{tcolorbox}
  \vspace{6pt}
  \begin{columns}[T]
    \begin{column}{0.55\textwidth}
      \hl{What they do:}
      \begin{itemize}
        \item Model optimal targeting in a static advertising market
        \item Firm chooses which consumer segments to target and how intensively
        \item Compare offline (coarse) vs.\ online (fine) targeting technologies
      \end{itemize}
    \end{column}
    \begin{column}{0.42\textwidth}
      \hl{Relevance:} Starting point for the dynamic targeting extension. This project
      extends to a dynamic setting where targeting depends on the consumer's $\bar{X}_{it}$
      state
    \end{column}
  \end{columns}
\end{frame}
\begin{frame}{Hosanagar, Fleder, Lee \& Buja (2014): Recommender Systems}
  \begin{tcolorbox}[colback=red!5!white, colframe=UGARed, title=\textbf{Citation}]
    Hosanagar, K., Fleder, D., Lee, D., and Buja, A. (2014). Will the global village
    fracture into tribes? Recommender systems and their effects on consumer fragmentation.
    \textit{Management Science}, 60(4), 805--823.
  \end{tcolorbox}
  \vspace{6pt}
  \begin{columns}[T]
    \begin{column}{0.55\textwidth}
      \hl{What they do:}
      \begin{itemize}
        \item Ask whether algorithmic recommendations fragment or homogenize consumption
        \item Construct attribute-distance measures across product pairs
        \item Find recommendations \emph{homogenize} consumption despite personalization
      \end{itemize}
    \end{column}
    \begin{column}{0.42\textwidth}
      \hl{Relevance:} Attribute distance measures are repurposed here on the \emph{firm}
      side to determine when and whom to advertise to; homogenization result informs
      expected equilibrium outcomes
    \end{column}
  \end{columns}
\end{frame}

\begin{frame}{Song (2025): Personalization Field Experiment}
  \begin{tcolorbox}[colback=red!5!white, colframe=UGARed, title=\textbf{Citation}]
    Song, Yu. (2025). Optimizing multi-stage personalization in the customer journey.
    \textit{Working Paper}.
  \end{tcolorbox}
  \vspace{6pt}
  \begin{columns}[T]
    \begin{column}{0.55\textwidth}
      \hl{What they do:}
      \begin{itemize}
        \item Partner with a major e-tailer for a randomized field experiment
        \item Vary personalization intensity across stages of the consumer search journey
        \item Identify causal effect of personalization on consumption patterns
      \end{itemize}
    \end{column}
    \begin{column}{0.42\textwidth}
      \hl{Relevance:} Practical implementation of Wei \& Jiang (2024) neural methods in
      a live commercial setting; provides external validity for the structural model's
      policy predictions
    \end{column}
  \end{columns}
\end{frame}

% ============================================================
\section{Summary}
% ============================================================

\begin{frame}{How the Papers Fit Together}
  \begin{columns}[T]
    \begin{column}{0.48\textwidth}
      \hl{Consumer side}
      \begin{itemize}\small
        \item Pollak (1970): habit formation microfoundations
        \item Dixit-Stiglitz (1977): love of variety
        \item McAllister-Pessemier (1982): behavioral evidence
        \item Becker-Murphy (1988): rational addiction stock
        \item Bronnenberg et al.\ (2012): empirical persistence
      \end{itemize}
    \end{column}
    \begin{column}{0.48\textwidth}
      \hl{Firm / advertising side}
      \begin{itemize}\small
        \item Nerlove-Arrow (1962): goodwill dynamics
        \item Erdem-Keane (1996): learning from ads
        \item Dube et al.\ (2005): ad timing and persistence
        \item Bergemann-Bonatti (2011): optimal targeting
        \item Hosanagar et al.\ (2014): fragmentation vs.\ homogenization
        \item Song (2025): field experiment, neural estimation
      \end{itemize}
    \end{column}
  \end{columns}
\end{frame}

\begin{frame}{Contribution of This Paper}
  \begin{tcolorbox}[colback=red!5!white, colframe=UGARed, title=\textbf{Where This Project Sits}]
    We connect the love-of-variety / habit-formation demand literature to the dynamic
    targeting literature by making $\bar{X}_{it}$ the \hl{sufficient statistic} for both
    consumer switching behavior and firm advertising decisions.
  \end{tcolorbox}
  \vspace{8pt}
  \begin{enumerate}
    \item \hl{Novel mechanism:} Utility depends on deviation from experienced mean ---
          nests Pollak and Dixit-Stiglitz in a single dynamic framework
    \vspace{4pt}
    \item \hl{Dynamic targeting:} Extends Bergemann-Bonatti to a setting where the firm's
          optimal ad policy is a function of the consumer's history
    \vspace{4pt}
    \item \hl{Estimation:} Bring the model to Nielsen scanner data using BLP/PyBLP
          with potential neural-network moments
  \end{enumerate}
\end{frame}

\end{document}
